% Part: normal-modal-logic
% Chapter: filtrations
% Section: euclidean-filtrations

\documentclass[../../../include/open-logic-section]{subfiles}

\begin{document}

\olfileid[hi]{nml}{fil}{euc}

\olsection{यूक्लिडीय प्रतिरूपों के निस्यंदन}

\olref[acc]{sec} का दृष्टिकोण उन प्रतिरूपों के लिए काम नहीं करता जो
यूक्लिडीय, या सीरियल और यूक्लिडीय हैं। \olref{fig:ser-eucl} के ऊपर वाले
प्रतिरूप पर विचार करें, जो यूक्लिडीय और सीरियल, दोनों है। $\Gamma = \{p,
\Box p \}$ रखें। $\Gamma$ के माध्यम से निस्यंदन लेते समय $[w_1] =
[w_3]$, क्योंकि $w_1$ और $w_3$ ही वे जगत हैं जो $\Gamma$ पर सहमत हैं।
प्रत्येक निस्यंदन में $\mModel{M}$ से उत्तराधिकार में मिला तीर भी होगा,
जैसा \olref{fig:ser-eucl2} में दर्शाया गया है। वह प्रतिरूप यूक्लिडीय नहीं
है। इसके अतिरिक्त, हम उसे यूक्लिडीय बनाने के लिए उसमें तीर नहीं जोड़ सकते।
हमें $[w_2]$ और $[w_4]$ के बीच दोनों दिशाओं में तीर जोड़ने होंगे, और फिर
$w_2$ और~$w_5$ के बीच भी। किंतु $\Box p$, $w_2$ पर सत्य होना चाहिए,
जबकि $p$, $w_5$ पर असत्य है।

\begin{figure}[htpb]
  \centering
  \begin{tikzpicture}[modal]
    \node[world] (w1)
          [label={left: \mFalse{p}},
            label={below: $\mSat{{}}{\Box p}$}] {$w_1$}; 
    \node[world] (w2)
          [label={right: \mTrue{p}},
            label = {below: $\mSat{{}}{\Box p}$},
            right=of w1] {$w_2$}; 
    \draw[->] (w1) to (w2);
    \node[world] (w3)
          [label={left: \mFalse{p}},
            label={below: $\mSat{{}}{\Box p}$},
            below=of w1] {$w_3$}; 
    \node[world] (w4)
          [label={right:\mTrue{p}},
            label={below: $\mSat/{{}}{\Box p}$},
            right=of w3] {$w_4$};
    \draw[->] (w3) to (w4);
    \draw[reflexive above] (w4) to (w4);     
    \node[world] (w5)
          [label={right: \mFalse{p}},
            label=below:$\mSat/{{}}{\Box p}$,
            right=of w4] {$w_5$}; 
    \draw[->, bend left] (w4) to (w5);
    \draw[->, bend left] (w5) to (w4);
    \draw[reflexive above] (w5) to (w5);     
  \end{tikzpicture}
  \caption{एक सीरियल और यूक्लिडीय प्रतिरूप।}
  \ollabel{fig:ser-eucl}
\end{figure}

\begin{figure}[ht]
  \centering
  \begin{tikzpicture}[modal,every text node part/.style={align=center}]
    \node[world] (w1)
          [label={left: \mFalse{p}},
            label={right:$[w_1]=[w_3]$},
            label={below: $\mSat{{}}{\Box p}$}] {$[w_1]$}; 
    \node[world] (w2)
         [label={right: \mTrue{p}},
           label = {below: $\mSat{{}}{\Box p}$},
           right=of w1, yshift=1.5cm] {$[w_2]$}; 
    \draw[->] (w1) to (w2);
    \node[world] (w4)
         [label={right:\mTrue{p}},
           label={below: $\mSat/{{}}{\Box p}$},
           right=of w1,yshift=-1.5cm] {$[w_4]$};
    \draw[->] (w1) to (w4);
    \draw[reflexive above] (w4) to (w4);     
    \node[world] (w5)
         [label={right: \mFalse{p}},
             label=below:$\mSat/{{}}{\Box p}$,
             right=of w4] {$[w_5]$}; 
    \draw[->, bend left] (w4) to (w5);
    \draw[->, bend left] (w5) to (w4);
    \draw[reflexive above] (w5) to (w5);     
  \end{tikzpicture}
  \caption{\olref{fig:ser-eucl} के प्रतिरूप का निस्यंदन।}
  \ollabel{fig:ser-eucl2}
\end{figure}

विशेषतः, यूक्लिडीय निस्यंदन पाने के लिए उप-!!{formula}ों के अधीन संवृत
मनमाने $\Gamma$ के माध्यम से निस्यंदनों पर विचार करना पर्याप्त नहीं है।
इसके स्थान पर हमें ऐसे समुच्चयों $\Gamma$ पर विचार करना होगा जो
\emph{मोडलतः संवृत} हैं (\olref[pre]{defn:modallyclosed} देखें)। वाक्यों
के ऐसे समुच्चय अपरिमित होते हैं, इसलिए वे तुरंत परिमित प्रतिरूप गुण या
संबद्ध तंत्र की निर्णेयता नहीं देते।

\begin{thm}\ollabel{thm:modal-closed-filt}
  मान लें कि $\Gamma$ मोडलतः संवृत है, $\mModel{M}=\tuple{W,R,V}$, और
  $\mModel{M^*} = \tuple{W^*,R^*,V^*}$, $\mModel{M}$ का स्थूलतम
  निस्यंदन है।
  \begin{enumerate}
  \item यदि $\mModel{M}$ सममित है, तो $\mModel{M^*}$ भी।
  \item यदि $\mModel{M}$ संक्रमणशील है, तो $\mModel{M^*}$ भी।
  \item यदि $\mModel{M}$ यूक्लिडीय है, तो $\mModel{M^*}$ भी।
  \end{enumerate}
\end{thm}

\begin{proof}
\begin{enumerate}
  \item यदि $\mModel{M^*}$ स्थूलतम निस्यंदन है, तो परिभाषा से
    $R^*[u][v]$ तभी सत्य है जब $C_1(u,v)$। संक्रमणशीलता के लिए $C_1(u,v)$
    और $C_1(v,w)$ मान लें; हमें $C_1(u,w)$ दिखाना है।
    \iftag{prvBox}{मान लें $\mSat{M}{\Box !A}[u]$; तब
      $\mSat{M}{\Box\Box!A}[u]$, क्योंकि \Ax{4} सभी संक्रमणशील
      प्रतिरूपों में मान्य है; संवृत्ति से $\Box\Box!A \in \Gamma$, अतः
      $C_1(u,v)$ से $\mSat{M}{\Box!A}[v]$ भी, और $C_1(v,w)$ से
      $\mSat{M}{!A}[w]$ भी।}{}
    \iftag{prvDiamond}{मान लें $\mSat{M}{!A}[w]$; तब $C_1(v,w)$ से
      $\mSat{M}{\Diamond !A}[v]$, क्योंकि मोडल संवृत्ति से $\Diamond !A
      \in \Gamma$। $C_1(u,v)$ से $\mSat{M}{\Diamond\Diamond !A}[u]$
      मिलता है, क्योंकि मोडल संवृत्ति से $\Diamond\Diamond!A \in
      \Gamma$। चूँकि $\Ax{4}_\Diamond$ सभी संक्रमणशील प्रतिरूपों में
      मान्य है, $\mSat{M}{\Diamond!A}[u]$।}{}
  \item अभ्यास। इस तथ्य का उपयोग कीजिए कि \Ax{5} और $\Ax{5_\Diamond}$,
    दोनों सभी यूक्लिडीय प्रतिरूपों में मान्य हैं।
  \item अभ्यास। इस तथ्य का उपयोग कीजिए कि \Ax{B} और $\Ax{B_\Diamond}$,
    सभी सममित प्रतिरूपों में मान्य हैं।
\end{enumerate}
\end{proof}

\begin{prob}
  \olref[nml][fil][euc]{thm:modal-closed-filt} का प्रमाण पूरा कीजिए।
\end{prob}

\end{document}
